\section{结论}

本文把自由形态傅里叶算子超表面的逆向设计重新表述为一个定义在波长、角度和偏振上的响应中心问题。在这一表述下，方法整体构成一条连贯的物理闭环：制造约束结构生成、仿真标注、任务打分、数据增强、前向代理训练、条件扩散采样以及基于仿真验证的重评分。本文最核心的观点是，光学目标应该在最终器件也要被验证的那个响应空间中被定义，这也使本文与目标驱动的光子逆向设计以及响应一致的电磁建模路线形成了更直接的联系 \cite{an2019deep,ma2019probabilistic,phillips2021ensemble}。

基于这一定义，一个兼具大NA、高透过率和宽带工作的二阶微分器在多个波长下都达到了领先水平的验证表现；同时，偏振无关、偏振复用、四阶微分、低通以及时空微分等目标也都可以在同一套统一的推理框架中实现具有竞争力的高水平结果。因而，本文更重要的贡献在于证明多个算子家族可以在同一套物理一致的设计语言里被描述、生成和评估。

更广义地看，自由形态超表面逆向设计可以被理解为在电磁验证约束下的目标条件可行集发现问题。本文给出了这一原则的一种具体实现，并自然指向了未来的扩展方向：更丰富的偏振控制、更完整的复响应目标、更广的物理设计空间、更严格的制造模型以及更大规模的 physics-guided 数据闭环。
